% Part: first-order-logic
% Chapter: model-theory
% Section: substructures

\documentclass[../../../include/open-logic-section]{subfiles}

\begin{document}

\olfileid{mod}{bas}{sub}
\olsection{Deel\printtoken{p}{structure} (substructuren)}

Het !!{domain} van !!a{structure}~$\Struct{M}$ kan binnen dat van een
andere structuur~$\Struct{M'}$ liggen.  Dat alleen maakt $\Struct{M}$
nog geen ``deel'' van $\Struct{M'}$.  Daarvoor moet niet alleen
$\Domain{M} \subseteq \Domain{M'}$ gelden.  Bovendien moeten
$\Struct{M}$ en $\Struct{M'}$ op het gedeelde deel~$\Domain{M}$ aan
de symbolen van de taal dezelfde betekenis geven.

\begin{defn}
\ollabel{defn:substructure}
Neem !!{structure}s $\Struct M$ en $\Struct M'$ voor dezelfde
taal~$\Lang L$.  We noemen $\Struct M$ een \emph{deel!!{structure}}
(ook: substructuur) van $\Struct M'$.  Omgekeerd noemen we $\Struct M'$
een \emph{uitbreiding} van $\Struct M$.  We schrijven dit als
$\Struct M \substruct \Struct M'$.  Dat geldt precies wanneer:
\begin{enumerate}
\item $\Domain{M} \subseteq \Domain{M'}$,
\item elke constante $c \in \Lang L$ in beide structuren dezelfde
  betekenis heeft: $\Assign{c}{M} =
    \Assign{c}{M'}$;
\item elke $n$-plaatsige !!{function} $f \in \Lang L$ bij alle
  invoeren uit het kleinere domein dezelfde uitkomst geeft:
  $\Assign{f}{M}(a_1, \dots, a_n) = \Assign{f}{M'}(a_1, \dots, a_n)$
  voor alle $a_1$, \dots, $a_n \in \Domain{M}$.
\item elk $n$-plaatsig !!{predicate} $R \in \Lang L$ voor precies
  dezelfde tupels uit het kleinere domein geldt: $\langle
  a_1, \dots, a_n\rangle \in \Assign{R}{M}$ iff $\langle a_1, \dots,
  a_n\rangle \in \Assign{R}{M'}$ voor alle $a_1$, \dots, $a_n \in
  \Domain{M}$.
\end{enumerate}
\end{defn}

\begin{rem}
\ollabel{rem:substructure}
Als de taal geen constanten of !!{function}s bevat, is het eenvoudiger.
Neem om het even welke $N \subseteq \Domain{M}$.  Daarmee krijg je een
deel!!{structure}~$\Struct{N}$ van $\Struct M$.  Het !!{domain} daarvan
is $\Domain{N} = N$.  De relaties leggen we vast met $\Assign{R}{N} =
\Assign{R}{M} \cap N^n$.
\end{rem}

% prove something about this? Examples?

\end{document}
